% Three basic quadrature rules on one interval, drawn.
%
%   figures/tikz/ni-quadrature-rules.tex
%       -> media/figures/ni-quadrature-rules.{png,svg}
%
% Lecture 6 (Numeric Integration), Quadrature. Added 2026-10-07 (format pass):
% the lecture had no picture of a quadrature rule. One concept: each rule
% replaces f by the polynomial through its nodes and integrates that polynomial.
%
% GEOMETRY (x unit 1.45 cm, y unit 1.15 cm). f(x) = 1 + 0.8 sin(1.6 x) + 0.3 x
% on [a, b] = [0, 2], midpoint m = 1.
%   f(0) = 1, f(1) = 2.09966, f(2) = 1.55330.
%   (1) midpoint:  constant f(1), area 4.199
%   (2) trapezoid: chord through f(0), f(2), area 2.553
%   (3) Simpson:   parabola through all three, p(x) = 1 + 1.92267 x - 0.82301 x^2,
%                  area (2/6)(f(0) + 4 f(1) + f(2)) = 3.651
%   exact integral 2 + 0.5 (1 - cos 3.2) + 0.6 = 3.599.
% The panels show the shape of each approximation; the numbers are not printed.
%
% GREYSCALE: f is solid black, the approximating polynomial is dashed, the
% approximate area is hatched, nodes are filled dots. Colour is redundant.
\definecolor{okabeblue}{HTML}{0072B2}

\begin{tikzpicture}[
    x=1.45cm, y=1.15cm,
    font=\small,
    lab/.style={font=\scriptsize, inner sep=1.5pt},
    fcurve/.style={very thick, black},
    approx/.style={okabeblue, thick, dashed},
    area/.style={pattern=north east lines, pattern color=okabeblue!70},
    node dot/.style={circle, fill=black, inner sep=0pt, minimum size=4.2pt},
  ]
  \def\fcurvepath{plot[domain=0:2, smooth, samples=60, variable=\t]
    ({\t},{1+0.8*sin(deg(1.6*\t))+0.3*\t})}
  \foreach \k/\name in {0/{(1) midpoint}, 1/{(2) trapezoid}, 2/{(3) Simpson}} {
    \begin{scope}[xshift=\k*4.2cm]
      % approximate area
      \ifcase\k
        \fill[area] (0,0) rectangle (2,2.09966);
        \draw[approx] (0,2.09966) -- (2,2.09966);
        \draw[approx] (0,0) -- (0,2.09966);
        \draw[approx] (2,0) -- (2,2.09966);
      \or
        \fill[area] (0,0) -- (0,1) -- (2,1.55330) -- (2,0) -- cycle;
        \draw[approx] (0,1) -- (2,1.55330);
        \draw[approx] (2,0) -- (2,1.55330);
      \or
        \fill[area] (0,0) -- plot[domain=0:2, smooth, samples=40, variable=\t]
          ({\t},{1+1.92267*\t-0.82301*\t*\t}) -- (2,0) -- cycle;
        \draw[approx] plot[domain=0:2, smooth, samples=40, variable=\t]
          ({\t},{1+1.92267*\t-0.82301*\t*\t});
        \draw[approx] (2,0) -- (2,1.55330);
      \fi
      % axes
      \draw[->, black!60] (-0.12,0) -- (2.3,0) node[right, lab] {$x$};
      \draw[->, black!60] (0,0) -- (0,2.6);
      % the integrand
      \draw[fcurve] \fcurvepath;
      \node[lab, anchor=south west] at (1.55,2.0) {$f$};
      % nodes
      \ifcase\k
        \node[node dot] at (1,2.09966) {};
      \or
        \node[node dot] at (0,1) {};
        \node[node dot] at (2,1.55330) {};
      \or
        \node[node dot] at (0,1) {};
        \node[node dot] at (1,2.09966) {};
        \node[node dot] at (2,1.55330) {};
      \fi
      % ticks
      \foreach \xx/\ll in {0/a, 1/m, 2/b} {
        \draw[black!60] (\xx,0) -- (\xx,-0.06);
        \node[lab, anchor=north] at (\xx,-0.06) {$\ll$};
      }
      \node[font=\small, anchor=north] at (1,-0.38) {\name};
    \end{scope}
  }
\end{tikzpicture}
